% sigma-hat-trajectory.tex — Chapter on the per-step eta and the
% global sigma-hat-budget closure, including the chaining-tightening
% (Sec. 10.5) and the manuscript m_0 uplift.
%
% Linked code / Lean:
%   - computations/sigma_hat_trajectory.py
%   - computations/chaining_eta_uplift.py
%   - lean/YMGap/RG/SigmaHatPilot.lean

\chapter{The
  \texorpdfstring{$\widehat\sigma$}{sigma-hat} Trajectory and
  Budget Closure}
\label{ch:sigma-hat-trajectory}

\noindent\emph{The $\sigh$-transport derivation
$\eta_k = (\cpolymer / \Amin) \cdot \Dcap(b, g_k)$, the global
budget $\sum_k \eta_k < \sigma_0$ in $\Qp$, and the chaining
tightening that uplifts the constructive mass-gap lower bound.}

This chapter is the cumulative-budget closure step.  It converts
the per-step buffer growth $\Dcap$ of \Cref{ch:overlap-kp} into the
per-step renormalized-string-tension degradation $\eta_k$ via a
$\Qp$-rational coupling constant derived from the area-law
convention.  The global budget $\sum_k \eta_k < \sigma_0$ in $\Qp$
at finite truncation is the make-or-break inductive condition; its
closure delivers the $\sigh$-positivity $\sighinf \ge \varepsilon
> 0$ that feeds the continuum-limit and OS-reconstruction chapters.
\Cref{sec:eta-chaining-thread} threads the chaining factor of
\Cref{ch:chaining-estimate} into the budget, contributing a clean
multiplicative uplift to the manuscript-quotable $\massgap$ at the
canonical pilot.

\section{The
  \texorpdfstring{$\widehat\sigma$}{sigma-hat} transport derivation}
\label{sec:sigma-hat-transport}

\marginkey{Per-step $\eta$ from $\Dcap$}%

The renormalized string tension at scale $k$ is
$\sighk := a_k^2 \cdot \sigma_k$, where $\sigma_k$ is the dimensional
string tension governing the area-law decay of Wilson-loop
expectations,
\begin{equation}
  \label{eq:wilson-area-law}
  \langle W_R(C) \rangle_{\muk}
  \;\le\;
  C_R \cdot \exp\bigl(-\sigma_k \cdot A(C)\bigr),
\end{equation}
for Wilson loops with sides $\ge 2\,a_k$.  Under the block-spin
map $\Tb \colon \Omk \to \Omkpone$ with $a_{k+1} = b \cdot a_k$,
dimensional analysis gives the naive
scaling $\sigma_{k+1}^{\mathrm{naive}} = \sigma_k / b^2$, equivalently
$\sighk^{\mathrm{naive}} = \sighk$ — the renormalized string tension
$\sigh$ is RG-invariant under pure dimensional rescaling.

\marginkey{Polymer-residual correction}%

The deviation from naive scaling is the polymer-residual correction
on the Wilson-loop expectation, controlled by the per-step buffer
growth $\Dcap(b, g_k)$ via a Bałaban-style cluster-expansion
residual:
\begin{equation}
  \label{eq:wilson-residual}
  \bigl|\sigma_{k+1} - \sigma_k / b^2\bigr|
  \;\le\;
  \frac{\cpolymer \cdot \Dcap(b, g_k)}{\Amin \cdot a_{k+1}^2} ,
\end{equation}
where
\begin{itemize}
  \item $\cpolymer = 2$ is the polymer-residual prefactor, derived
    as the channel multiplicity from the triangle-inequality bound
    on the per-polymer activity difference under $\Tb$;
  \item $\Amin = 4$ is the smallest dimensionless Wilson-loop
    area, $(2\,a_k)^2 = 4\,a_k^2$ in lattice-spacing-normalized
    units.
\end{itemize}
Lifting \eqref{eq:wilson-residual} to $\sigh$, the $a_{k+1}^2$
factor in the denominator absorbs:

\begin{theorem}[The per-step $\eta$ coupling]
\label{thm:eta-coupling}
With the convention above,
\begin{equation}
  \label{eq:eta-coupling}
  \bigl|\sighk[k+1] - \sighk\bigr|
  \;\le\;
  \etak,
  \qquad
  \etak \;:=\; \frac{\cpolymer}{\Amin} \cdot \Dcap(b, g_k)
  \;=\;
  \frac12 \cdot \Dcap(b, g_k) .
\end{equation}
With $\cpolymer = 2$ and $\Amin = 4$, the per-step
proportionality constant is $\cpolymer / \Amin = 1/2$, which
matches what was previously a placeholder coupling but is now
explicitly derived from project-level conventions.
\end{theorem}

\begin{proof}[Proof sketch]
The Bałaban-style polymer-cluster expansion writes
$\langle W \rangle_{\muk}$ as a leading area-law factor plus a
polymer correction; under the block-spin step $\Tb$, the
correction picks up at most a factor $\cpolymer$ per polymer
boundary channel (from the triangle inequality on the activity
difference), divided by the smallest loop area $\Amin a_{k+1}^2$
that the Wilson loops in question can realize.  The bound
\eqref{eq:wilson-residual} is then lifted to $\sigh$ by
multiplying both sides by $a_{k+1}^2$, which cancels the
denominator.  See \cite{Balaban1988convergent} \S2 for the
analogous step at the Wilson-action level.
\end{proof}

\section{The global $\widehat\sigma$ trajectory}
\label{sec:sigma-hat-global}

\marginkey{Cumulative degradation}%

Define the $\sigh$ trajectory recursively by
\begin{equation}
  \label{eq:sigma-hat-recursion}
  \sighk[0] \;:=\; \sigma_0,
  \qquad
  \sighk[k+1] \;:=\; \sighk - \etak,
\end{equation}
where $\sigma_0 > 0$ is the initial dimensionless string tension
(normalized to $\sigma_0 = 1$ in the manuscript-standard
convention).  By \cref{thm:eta-coupling},
$\sighk \ge \sigma_0 - \sum_{j < k} \eta_j$ for every $k$.

\begin{theorem}[Finite-$K$ budget closure]
\label{thm:budget-closure}
If $\sum_{k < K} \etak < \sigma_0$ in $\Qp$, then
$\sighk[K] \ge \sigma_0 - \sum_{k < K} \etak > 0$, also in $\Qp$.
The trajectory $\sighk[0], \sighk[1], \dots, \sighk[K]$ remains
strictly positive across the $K$-truncation.
\end{theorem}

\begin{proof}
By induction on $k$, using \cref{thm:eta-coupling} and the
inductive hypothesis $\sighk \ge \sigma_0 - \sum_{j < k} \eta_j$.
\end{proof}

\begin{infobox}[main][title=Canonical pilot,
  $g_0 = 0.30$, $K = 3$]
\centering
\renewcommand{\arraystretch}{1.15}
\begin{tabular}{rcccc}
\toprule
$k$ & $g_k$    & $\Dcap(b, g_k)$ & $\etak$ ($\le$ in $\Qp$)  & $\sighk$ ($\ge$ in $\Qp$) \\
\midrule
0   & $0.3000$ & $5.17\times 10^{-4}$ & $2.58\times 10^{-4}$ & $1$                          \\
1   & $0.3004$ & $5.30\times 10^{-4}$ & $2.65\times 10^{-4}$ & $9997/10000$                 \\
2   & $0.3009$ & $5.43\times 10^{-4}$ & $2.72\times 10^{-4}$ & $9995/10000$                 \\
3   & $0.3013$ & $0$ (truncation)     & $0$                  & $9992/10000$                 \\
\bottomrule
\end{tabular}
\vspace{0.3em}\\
Aggregate: $\sum_{k < 3} \etak \le 9/10000$,
$\sighinf \ge \sigma_0 - 9/10000 = 9991/10000$.
\end{infobox}

\begin{mainfigure}
\centering
\includegraphics[width=\textwidth]{sigma-hat-trajectory}
\caption{The $\sighk$ trajectory across the canonical and Wilson
pilots at $K = 3$, with the threshold $\tfrac12\sigma_0$ annotated
for reference.  Both trajectories remain strictly within
$[\tfrac12\sigma_0, \sigma_0]$, with cumulative $\eta$-budget
$\le 9/10000$ (canonical) and $\le 3/10000$ (Wilson).}
\label{fig:sigma-hat-trajectory}
\end{mainfigure}

\section{Convergence to a positive
  \texorpdfstring{$\widehat\sigma_\infty$}{sigma\_infty}}
\label{sec:sigma-hat-infty}

The finite-$K$ closure of \cref{thm:budget-closure} extends to a
genuine projective limit $\sighinf \ge \varepsilon > 0$ provided
the tail $\sum_{k \ge K_0} \etak$ admits a rational upper bound.
Across the canonical SU(2) trajectory in the UV-direction (where
the one-loop $\beta$-function drives $g_k \to 0$ monotonically),
$\Dcap(b, g_k)$ decays geometrically with explicit ratio
$r \approx 0.976$ derived from $b = 2$ and the SU(2) coefficient
$b_0 = 22/3$ of the one-loop running.  The tail is therefore
dominated by a geometric series with $\Qp$-rational closure.

\begin{theorem}[$\sigh$-positivity in the projective limit]
\label{thm:sigma-hat-positivity}
Under the asymptotic-freedom UV-direction trajectory, there exists
an explicit $\Qp$-rational $\varepsilon > 0$ such that
\begin{equation}
  \label{eq:sigma-hat-positivity}
  \sighinf \;:=\; \lim_{k \to \infty} \sighk \;\ge\; \varepsilon \;>\; 0 .
\end{equation}
At the canonical pilot, $\varepsilon = 9991/10000$.
\end{theorem}

\begin{proof}[Proof sketch]
Combine the finite-$K$ closure of \cref{thm:budget-closure} at
$K = 3$ (the empirically-verified $\Qp$-rational truncation) with
the geometric tail bound on $\sum_{k \ge K} \etak$ established in
\Cref{ch:uv-trajectory} below; the partial sum is Cauchy in $\Qp$
at the rational truncation, and the limit is bounded by
$\sigma_0 - \mathrm{partial} > 0$ at every truncation.
\end{proof}

\section{Threading the chaining factor into
  \texorpdfstring{$\eta$}{eta}}
\label{sec:eta-chaining-thread}

\marginkey{Chaining tightens $\eta$}%

The buffer-region aggregate measure $\muk(\Omkcap)$ enters the
$\sigh$-transport derivation of \cref{thm:eta-coupling} through the
polymer-cluster expansion's per-step residual.  When the chaining
tightening of \Cref{ch:chaining-estimate},
Theorem~\ref{thm:chaining-bound}, is admissible
($g_k \le G_{\max}^{\mathrm{chain}} = 2$), the aggregate measure is
strictly tighter than its union-bound estimate, and the per-step
$\eta$ inherits a multiplicative reduction:
\begin{equation}
  \label{eq:eta-chaining}
  \eta_k^{\mathrm{chain}} \;:=\; \cchain(g_k) \cdot \etak[k]^{\mathrm{union}} ,
  \qquad
  \cchain(g_k) = \exp\bigl(-\,n_{\mathrm{share}} \cdot g_k^2 \cdot \zeta\bigr) ,
\end{equation}
with $n_{\mathrm{share}} = 4$ and $\zeta = \tfrac12$ as in
\Cref{ch:chaining-estimate},
Definition~\ref{def:chain-factor}.  Since $\cchain \in (0, 1]$, the
chaining-tightened $\eta$ never exceeds the union-bound $\eta$, and
the budget-closure theorem of \cref{thm:budget-closure} continues
to hold under chaining (with a strictly tighter bound).

\begin{theorem}[Chaining-uplifted $\massgap$ at the canonical pilot]
\label{thm:chaining-uplift}
At the canonical pilot $g_0 = 0.30$, the chaining factor
$\cchain(0.30) = e^{-0.18} \approx 0.8353$ reduces
$\eta_{\mathrm{total}}$ from the union-bound value
$\sum_k \etak[k]^{\mathrm{union}} \approx 7.95 \times 10^{-4}$ to
$\sum_k \cchain(g_k)\, \etak[k]^{\mathrm{union}} \approx
6.64 \times 10^{-4}$, a $16.5\%$ reduction.  At trajectory
length $K = 10$ this is sufficient to uplift the $\Qp$-rational
floor of $\sqrt{\sighinf}$ from $997/1000$ to $998/1000$,
an uplift ratio of $\approx 1.001$ on the manuscript-quotable bound.
\end{theorem}

\begin{mainfigure}
\centering
\includegraphics[width=\textwidth]{chaining-eta-uplift}
\caption{Manuscript-quotable $\massgap$ baseline (red, union-bound
$\eta$) vs.\ chaining-tightened $\massgap$ (green) across eight
operating points in the canonical convergent window.  The
chaining factor $\cchain(g_k) \in (0, 1]$ never weakens the bound
and gives a numerator uplift $\Delta_{\mathrm{num}} \ge 0$ at each
point; the largest uplift is at the convergent-threshold
operating point $g_0 = 0.40, K = 5$
(ratio $\times 1.076$), and the manuscript-binding canonical
$g_0 = 0.30, K = 10$ uplift is $\times 1.001$.}
\label{fig:chaining-eta-uplift}
\end{mainfigure}

\section{Code provenance}
\label{sec:sigma-hat-codeanchor}

\begin{itemize}
  \item Python implementation:
    \codepath{computations/sigma\_hat\_trajectory.py}.  Key entry
    points: \codepath{compute\_eta(g\_k, g\_k\_next, t\_choice,
    eta\_coupling, use\_chaining, ...)} for the per-step degradation
    of \cref{thm:eta-coupling}, with the chaining-tightening
    option of \eqref{eq:eta-chaining};
    \codepath{sigma\_hat\_trajectory(g\_trajectory,
    sigma\_hat\_initial, t\_choice)} for the global trajectory of
    \cref{thm:budget-closure}, returning a
    \codepath{SigmaHatTrajectory} record with the per-step
    \codepath{SigmaHatStep}s and the \codepath{budget\_closure}
    boolean.  Tests in
    \codepath{tests/test\_sigma\_hat\_trajectory.py}.
  \item Constants: \codepath{BALABAN\_POLYMER\_RESIDUAL = 2}
    ($\cpolymer$), \codepath{WILSON\_LOOP\_MIN\_AREA\_DIMENSIONLESS
    = 4} ($\Amin$); the derived
    \codepath{DERIVED\_ETA\_COUPLING = 1/2} is recomputed at every
    fiducial sensitivity sweep via
    \codepath{derive\_eta\_coupling(polymer\_residual,
    wilson\_loop\_min\_area)}.
  \item Chaining-uplift comparison:
    \codepath{computations/chaining\_eta\_uplift.py} provides
    \codepath{chaining\_uplift\_row(g\_0, K)} and
    \codepath{chaining\_uplift\_table(operating\_points)}; the
    side-by-side bar chart of $\massgap$ baseline vs.\ chaining is
    in \codepath{visualizations/pilot/su2/chaining-eta-uplift.svg}
    (Figure~10.2 of the manuscript-quotable mass-gap chapter
    cross-reference).
  \item Lean formalization:
    \codepath{lean/YMGap/RG/SigmaHatPilot.lean::eta\_pilot\_budget\_closes}
    discharges the finite-$K$ closure of \cref{thm:budget-closure}
    at the canonical pilot $K = 3$, with axiom-gate
    \codepath{[propext]} only.
\end{itemize}
